\clearpage
\VerdierPhysicalPageAnchor{0242}%
\begingroup\small\setlength{\parskip}{0pt}\setlength{\emergencystretch}{1.5em}
\setlength{\abovedisplayskip}{4pt plus 2pt minus 2pt}%
\setlength{\belowdisplayskip}{4pt plus 2pt minus 2pt}%

\VerdierSemanticAnchor{ch4d-4-5-3-b-cont}
\begin{enumerate}
\item[1)] \textit{$X$ は $D^b(\mathcal A)$ の対象である。}
\item[2)] \textit{$X$ は $D^+(\mathcal A)$ の対象であり、$F$ は右定常である。}
\item[3)] \textit{$X$ は $D^-(\mathcal A)$ の対象であり、$F$ は左定常である。}
\item[4)] \textit{関手 $F$ は定常である。}
\end{enumerate}

\VerdierSemanticAnchor{ch4d-4-5-3-proof}
この命題は (4.5.2) と定義（第 II 章、4.4.2）から直ちに従う。

\VerdierSemanticAnchor{ch4d-4-5-4}
\medskip\noindent\textbf{定理 4.5.4.}\quad
\textit{$F_1,F_2\colon D^*(\mathcal A)\longrightarrow\mathcal B$
（$*=+,-,b,{}$ ``空''）を二つのコホモロジー的関手（それぞれ
$G_1,G_2\colon D^*(\mathcal A)\longrightarrow D^{*'}(\mathcal B)$ を
二つの完全関手）とし、$m\colon F_1\longrightarrow F_2$（それぞれ
$m\colon G_1\longrightarrow G_2$）を関手の射（それぞれ完全関手の射
（第 I 章、1.6.6.4））とする。各整数 $p$ に対して $m(X[p])$ が同型と
なるような $D^*(\mathcal A)$ の対象 $X$ によって定まる
$D^*(\mathcal A)$ の充満部分圏は、$D^*(\mathcal A)$ の厳密充満かつ
飽和な三角部分圏（第 II 章、2.1.6）であり、$\operatorname{Is}(m)$ と書く。}

\VerdierSemanticAnchor{ch4d-4-5-4-a}
\noindent\textit{a) $\mathcal A$ のすべての対象が
$\operatorname{Is}(m)$ の対象ならば、
$D^b(\mathcal A)\subset\operatorname{Is}(m)$ である。}

\VerdierSemanticAnchor{ch4d-4-5-4-b}
\noindent\textit{b) $\mathcal A$ のすべての対象が
$\operatorname{Is}(m)$ の対象であり、関手 $F_1$ と $F_2$（それぞれ
$G_1$ と $G_2$）が右定常ならば、
$D^+(\mathcal A)\cap D^*(\mathcal A)\subset\operatorname{Is}(m)$ である。}

\VerdierSemanticAnchor{ch4d-4-5-4-c}
\noindent\textit{c) $\mathcal A$ のすべての対象が
$\operatorname{Is}(m)$ の対象であり、関手 $F_1$ と $F_2$（それぞれ
$G_1$ と $G_2$）が左定常ならば、
$D^-(\mathcal A)\cap D^*(\mathcal A)\subset\operatorname{Is}(m)$ である。}

\VerdierSemanticAnchor{ch4d-4-5-4-d}
\noindent\textit{d) $\mathcal A$ のすべての対象が
$\operatorname{Is}(m)$ の対象であり、関手 $F_1$ と $F_2$（それぞれ
$G_1$ と $G_2$）が定常ならば、
$D^*(\mathcal A)=\operatorname{Is}(m)$ である。}

\VerdierSemanticAnchor{ch4d-4-5-4-e}
\noindent\textit{e) $\operatorname{Ob}(\mathcal A)$ の部分集合 $M$ が
存在して、次を満たすとする：}

\noindent\textit{$\alpha$) $\mathcal A$ のすべての対象は、$M$ に属する
ある対象の商対象と同型である；}

\noindent\textit{$\beta$) $M$ に属するすべての対象は
$\operatorname{Is}(m)$ の対象である；}

\noindent\textit{さらに関手 $F_1$ と $F_2$（それぞれ $G_1$ と $G_2$）が
左定常ならば、
$D^-(\mathcal A)\cap D^*(\mathcal A)\subset\operatorname{Is}(m)$ である。}

\VerdierSemanticAnchor{ch4d-4-5-4-f}
\noindent\textit{f) $\operatorname{Ob}(\mathcal A)$ の部分集合 $M$ が
存在して、次を満たすとする：}

\noindent\textit{$\alpha$) $\mathcal A$ のすべての対象は、$M$ に属する
ある対象の部分対象と同型である；}

\noindent\textit{$\beta$) $M$ に属するすべての対象は
$\operatorname{Is}(m)$ の対象である；}

\noindent\textit{さらに関手 $F_1$ と $F_2$（それぞれ $G_1$ と $G_2$）が
右定常ならば、
$D^+(\mathcal A)\cap D^*(\mathcal A)\subset\operatorname{Is}(m)$ である。}

\VerdierSemanticAnchor{ch4d-4-5-4-proof-start}
まず、定理 (4.5.4) の括弧内の主張は、関手 $G_1$ と $G_2$ を標準
コホモロジー的関手 $D^{*'}(\mathcal B)\longrightarrow\mathcal B$ と
合成することにより、括弧外の主張から従うことに注意する。a) を示す。
$X$ を
\endgroup
